\documentclass[10pt,a4paper]{amsart}
\usepackage[margin=19mm]{geometry}
\usepackage{amsmath,amssymb,mathtools}
\usepackage[scr=rsfs,cal=euler]{mathalfa}
\usepackage{xfrac}
\usepackage[unicode,hidelinks,bookmarks=false]{hyperref}
\newcommand{\ie}{i.e.\ }
\newcommand{\cf}{cf.\ }
\newcommand{\resp}{respectively\ }
\newcommand{\Subsection}{\subsection}
\providecommand{\nobreakdash}{\nobreak\hbox{-}}
\setlength{\emergencystretch}{2em}
\multlinegap0pt
\usepackage[shortlabels]{enumitem}
\usepackage{amssymb}
\usepackage{mathtools}
\usepackage{xcolor}
\usepackage{algorithm}\usepackage{algpseudocode}
\newtheorem{lemma}{Lemma}
\newtheorem{theorem}{Theorem}
\newcommand{\E}{\mathbb{E}}
\newcommand{\Lone}{L^1}
\newcommand{\Ppi}{\mathcal{P}_\pi}
\newcommand{\R}{\mathbb{R}}
\newcommand{\calF}{\mathcal{F}}
\newcommand{\calX}{\mathcal{X}}
\title{The Geometry of Knowing: From Possibilistic Ignorance to Probabilistic Certainty\\[6pt] \large A Measure-Theoretic Framework for Epistemic Convergence}
\author{Moriba Kemessia Jah, Ph.D.}
\date{Source: arXiv:2604.09614v1 (CC BY 4.0)}
\begin{document}
\maketitle

\noindent\emph{Source and licence.} The Geometry of Knowing: From Possibilistic Ignorance to Probabilistic Certainty\\[6pt] \large A Measure-Theoretic Framework for Epistemic Convergence. Version arXiv:2604.09614v1 (CC BY 4.0), distributed by arXiv under the Creative Commons Attribution 4.0 International licence, http://creativecommons.org/licenses/by/4.0/, as recorded in the arXiv licence metadata for this exact version and shown on the abstract page https://arxiv.org/abs/2604.09614v1. Full licence notice: \texttt{licenses/arxiv-2604.09614-CC-BY-4.0.txt}.\par\smallskip

\noindent\textit{Editorial scope.} Counted unit: the article's theorem thm:choquet\_lebesgue (source0.tex lines 468--494). The article's complete original proof of it is delivered with it (source0.tex lines 496--554, one \texttt{proof} environment of 59 lines). No mathematics is written, repaired or re-derived: the statement and the proof are the article's own lines, in the article's own notation, and no retained line is substituted or re-flowed.  Retained uncounted as the source-local context the proof needs: lines 253--258; lines 394--407.  Nothing was omitted from any retained span, so the package carries no omission record.  No retained line contains a citation, so the wrapper carries no bibliography and no bibliography entry.  No retained line contains a figure environment, an image-inclusion command or drawing code, so no image is delivered and the package declares no asset dependency.  Editorial formatting only, in addition: an AMS \texttt{amsart} preamble that restates the article's own declarations and macros used by the retained lines, namely the environments \texttt{lemma}, \texttt{theorem} on separate counters, together with the standard packages those lines need.  The retained environments print in this document as Lemma 1, Theorem 1 in order of appearance, whereas the article prints the same items with the numbering of its own full document.  The complete native source of the article as distributed in the arXiv e-print for this exact version is bundled unchanged as \texttt{sources/198219/source0.tex}.
\begin{equation}
\underline{\E}_{\Ppi}[f]
\;\le\; \int f\,d\mathrm{Ch}_\pi
\;\le\; \overline{\E}_{\Ppi}[f],
\label{eq:choquet_sandwich}
\end{equation}

\begin{lemma}[Expectation oscillation bound]
\label{lem:oscillation}
Let $\mathcal{P}$ be a credal set on $(\calX,\calF)$ and let
$f:\calX\to\R$ be bounded measurable with $\|f\|_\infty\le M$.
Define the total-variation diameter
$d_{\mathrm{TV}}(\mathcal{P}) \equiv \sup_{P,Q\in\mathcal{P}}\|P-Q\|_{\mathrm{TV}}
= \sup_{P,Q\in\mathcal{P}}\sup_{A\in\calF}|P(A)-Q(A)|$.
Then
\begin{equation}
\overline{\E}_\mathcal{P}[f] - \underline{\E}_\mathcal{P}[f]
\;\le\; 2M\cdot d_{\mathrm{TV}}(\mathcal{P}).
\label{eq:oscillation}
\end{equation}
\end{lemma}

\begin{theorem}[Choquet-to-Lebesgue limit]
\label{thm:choquet_lebesgue}
Let $\{(\calX,\calF,\mu)\}$ be a $\sigma$-finite measurable space.
Let $\{\pi_t\}_{t\ge 0}$ be a sequence of normalized, upper-semi\-continuous,
\emph{consonant} possibility distributions on $(\calX,\calF,\mu)$ with
associated credal sets $\{\mathcal{P}_{\pi_t}\}$.
Suppose:
\begin{enumerate}[label=\textup{(A\arabic*)},ref=A\arabic*]
\item\label{A1} \emph{Consonance}: each $\pi_t$ is consonant, so that
  inequality~\eqref{eq:choquet_sandwich} holds with $\Pi_t$ a
  $\infty$-alternating capacity.
\item\label{A2} \emph{Credal contraction}: $|\mathcal{P}_{\pi_t}|\to 1$,
  i.e., $\sup_{A\in\calF}[\Pi_t(A)-N_t(A)]\to 0$ as $t\to\infty$.
\item\label{A3} \emph{$L^1$ convergence}: $\pi_t\to p^*$ in $\Lone(\mu)$ for
  some $p^*\in\Lone_+(\mu)$ with $\int_\calX p^*\,d\mu=1$.
\item\label{A4} \emph{Domination}: there exists $g\in\Lone(\mu)$ with $g\ge 0$
  and $\pi_t(x)\le g(x)$ $\mu$-a.e.\ for all $t$.
\end{enumerate}
Then for every bounded measurable $f:\calX\to\R$,
\begin{equation}
\int f\,d\mathrm{Ch}_{\pi_t}
\;\longrightarrow\;
\int_\calX f(x)\,p^*(x)\,d\mu(x)
\quad\text{as }t\to\infty.
\label{eq:choquet_limit}
\end{equation}
\end{theorem}

\begin{proof}
We proceed in four steps.

\smallskip\noindent
\textbf{Step 1: Sandwich from consonance.}
By assumption~\eqref{A1}, each $\Pi_t$ is a $\infty$-alternating capacity,
so inequality~\eqref{eq:choquet_sandwich} gives
\begin{equation}
\underline{\E}_t[f]
\;\le\;
\int f\,d\mathrm{Ch}_{\pi_t}
\;\le\;
\overline{\E}_t[f],
\label{eq:sandwich_t}
\end{equation}
where $\underline{\E}_t[f]=\inf_{P\in\mathcal{P}_{\pi_t}}\int f\,dP$ and
$\overline{\E}_t[f]=\sup_{P\in\mathcal{P}_{\pi_t}}\int f\,dP$.

\smallskip\noindent
\textbf{Step 2: Credal contraction forces both bounds to $\int f\,dP^*$.}
By assumption~\eqref{A2}, $\sup_{A\in\calF}[\Pi_t(A)-N_t(A)]\to 0$.
Since $d_{\mathrm{TV}}(\mathcal{P}_{\pi_t})\le\sup_A[\Pi_t(A)-N_t(A)]$,
Lemma~\ref{lem:oscillation} gives
\begin{equation}
\bigl|\overline{\E}_t[f]-\underline{\E}_t[f]\bigr|
\;\le\; 2\|f\|_\infty\cdot d_{\mathrm{TV}}(\mathcal{P}_{\pi_t})
\;\le\; 2M\cdot\sup_{A\in\calF}[\Pi_t(A)-N_t(A)]
\;\to\; 0.
\label{eq:squeeze}
\end{equation}
Hence the upper and lower expectations squeeze together.
Since $|\mathcal{P}_{\pi_t}|\to 1$, both converge to $\int f\,dP^*$.

\smallskip\noindent
\textbf{Step 3: Identifying $\int f\,dP^*$ with the Lebesgue integral.}
By assumption~\eqref{A3}, $\pi_t\to p^*$ in $\Lone(\mu)$.
Given~\eqref{A4}, the dominated convergence theorem gives
\begin{equation}
\int_\calX f(x)\,\pi_t(x)\,d\mu(x)
\;\to\;
\int_\calX f(x)\,p^*(x)\,d\mu(x).
\label{eq:dct}
\end{equation}
Since the $\Lone$ limit $p^*$ is a probability density and $P^*$ is the unique
member of the collapsed credal set, $\int f\,dP^*=\int_\calX f\,p^*\,d\mu$.

\smallskip\noindent
\textbf{Step 4: Combining.}
By~\eqref{eq:sandwich_t}, \eqref{eq:squeeze}, and \eqref{eq:dct},
\[
\Bigl|\int f\,d\mathrm{Ch}_{\pi_t}
  - \int_\calX f\,p^*\,d\mu\Bigr|
\;\le\;
\bigl|\overline{\E}_t[f]-\underline{\E}_t[f]\bigr|
+ \Bigl|\underline{\E}_t[f] - \int f\,dP^*\Bigr|
\;\to\; 0.
\]
This completes the proof.
\end{proof}

\end{document}
